\documentclass[10pt,a4paper]{amsart}
\usepackage[margin=19mm]{geometry}
\usepackage{amsmath,amssymb,mathtools}
\usepackage[scr=rsfs,cal=euler]{mathalfa}
\usepackage{xfrac}
\usepackage[unicode,hidelinks,bookmarks=false]{hyperref}
\newcommand{\ie}{i.e.\ }
\newcommand{\cf}{cf.\ }
\newcommand{\resp}{respectively\ }
\newcommand{\Subsection}{\subsection}
\providecommand{\nobreakdash}{\nobreak\hbox{-}}
\setlength{\emergencystretch}{2em}
\multlinegap0pt
\usepackage{amssymb}
\usepackage{enumerate}
\usepackage{bm}
\usepackage{float}
\usepackage{xcolor}
\usepackage{tikz}
\newtheorem{thm}{Theorem}
\title{On $S$-prime and $S$-primary elements in multiplicative lattices}
\author{Sachin Sarode*, Chetan Patil** and Vinayak Joshi***}
\date{Source: arXiv:2606.11932v1 (CC BY 4.0)}
\begin{document}
\maketitle

\noindent\emph{Source and licence.} On $S$-prime and $S$-primary elements in multiplicative lattices. Version arXiv:2606.11932v1 (CC BY 4.0), distributed by arXiv under the Creative Commons Attribution 4.0 International licence, http://creativecommons.org/licenses/by/4.0/, as recorded in the arXiv licence metadata for this exact version and shown on the abstract page https://arxiv.org/abs/2606.11932v1. Full licence notice: \texttt{licenses/arxiv-2606.11932-CC-BY-4.0.txt}.\par\smallskip

\noindent\textit{Editorial scope.} Counted unit: the article's thm thm (source0.tex lines 443--448). The article's complete original proof of it is delivered with it (source0.tex lines 450--464, one \texttt{proof} environment of 15 lines). No mathematics is written, repaired or re-derived: the statement and the proof are the article's own lines, in the article's own notation, and no retained line is substituted or re-flowed.  No further source-local context is needed: every cross-reference of every retained line resolves inside the two delivered spans.  Nothing was omitted from any retained span, so the package carries no omission record.  No retained line contains a citation, so the wrapper carries no bibliography and no bibliography entry.  No retained line contains a figure environment, an image-inclusion command or drawing code, so no image is delivered and the package declares no asset dependency.  Editorial formatting only, in addition: an AMS \texttt{amsart} preamble that restates the article's own declarations and macros used by the retained lines, namely the environments \texttt{thm} on separate counters, together with the standard packages those lines need.  The retained environments print in this document as Theorem 1 in order of appearance, whereas the article prints the same items with the numbering of its own full document.  The complete native source of the article as distributed in the arXiv e-print for this exact version is bundled unchanged as \texttt{sources/202585/source0.tex}.
\begin{thm}%\label{T4.6}
	Let $R$ be a commutative ring with identity and $S$ a multiplicative subset of $R$. 
	Let $L=Id(R)$ denote the multiplicative lattice of all ideals of $R$, and define 
	$S_L=\{(a)\in L \mid a\in S\}.$
	Then an ideal $P$ of $R$ is a weakly $S$-prime ideal of $R$ if and only if $P$ (viewed as an element of $L$) is a weakly $S_L$-prime element of $L$.
\end{thm}

\begin{proof}
	Suppose $P$ is a weakly $S$-prime ideal of $R$. Then $S\cap P=\varnothing$.  We claim that $t\not\leqq P$ for all $t\in S_L$. Suppose on the contrary that  there exists $t=(a)\in S_L$ with $t\le P$. But then $a\in P\cap S$, a contradiction.  
	Hence $t\not\leqq P$ for all $t\in S_L$.  
	
	Now, let $ I,J\in L$ with $0\neq IJ \leq P$, i.e., $\{0\}\neq IJ\subseteq P$. Since $P$ is weakly $S$-prime, there exists $s\in S$ such that $sI\subseteq P$ or $sJ\subseteq P$.  
	Equivalently, $(s)I\leq  P$ or $(s)J\leq P$ in $L$.  
	Thus $P$ is a weakly $S_L$-prime element of $L$.
	
	Conversely, assume that $P$ is a weakly $S_L$-prime element of $L$.  
	By the definition, $t\not\leq P$ for every $t\in S_L$. If some $a\in S\cap P$, then $(a)\in S_L$ with $(a)\le P$, a contradiction. Hence $S\cap P=\varnothing$.  
	
	Further, since $P$ is weakly $S_L$-prime, there exists $(s)\in S_L$ (for some $s\in S$) such that for all $I,J\in L$,
	$0\neq IJ\leq  P$ implies that $ (s)I\leq P \ \text{ or }\ (s)J\leq P$.
	That is, $sI\subseteq P$ or $sJ\subseteq P$. Thus $P$ is a weakly $S$-prime ideal of $R$.
\end{proof}

\end{document}
